% Chapter 1: Introduction and Research Objectives
% Deep Learning for Archaeological Site Detection from LiDAR-Derived Imagery
% A Multi-Version YOLO Consensus Approach

\documentclass[12pt,a4paper]{article}

% ── Packages ──────────────────────────────────────────────────
\usepackage[utf8]{inputenc}
\usepackage[T1]{fontenc}
\usepackage{newtxtext}
\usepackage{newtxmath}
\usepackage[margin=2.5cm]{geometry}
\usepackage{setspace}
\onehalfspacing

\usepackage[numbers,sort&compress]{natbib}
\usepackage{graphicx}
\usepackage{hyperref}
\usepackage{abstract}
\usepackage{titlesec}
\usepackage{enumitem}

% ── Metadata ──────────────────────────────────────────────────
\title{\textbf{Deep Learning for Archaeological Site Detection from LiDAR-Derived Imagery: \\
A Multi-Version YOLO Consensus Approach}\\[6pt]
\large Chapter 1: Introduction and Research Objectives}

\author{}
\date{June 2026}

\begin{document}
\maketitle

% ======================================================================
\section{Introduction}
% ======================================================================

The detection and documentation of archaeological sites at landscape scale remains one of the
central challenges in archaeological prospection. Across Europe, and particularly in the upland
regions of Scotland, thousands of archaeological features---ranging from prehistoric roundhouses
and shielings to burial cairns and field systems---lie unrecorded beneath vegetation, peat, and
modern land use \citep{banaszek2018only}. Systematic ground survey, while accurate, is
prohibitively slow and expensive when applied across entire regions. The development of airborne
remote sensing technologies over the past two decades has fundamentally altered this landscape,
offering the first realistic prospect of comprehensive, large-scale archaeological mapping.

Airborne Light Detection and Ranging (LiDAR) has emerged as the most transformative remote
sensing technology for archaeological prospection in vegetated terrain. By emitting laser pulses
that penetrate vegetation canopies and recording the reflected signal, LiDAR produces
high-resolution digital terrain models (DTMs) that reveal micro-topographic features invisible to
aerial photography or satellite imagery \citep{crutchley2010light}. When processed through
specialised visualisation techniques---including the Local Relief Model \citep{hesse2010lidar},
Sky-View Factor \citep{kokalj2011application}, and multiscale hillshading---LiDAR-derived
imagery can expose subtle archaeological earthworks with remarkable clarity. These
visualisations have led to the discovery of tens of thousands of previously unknown sites across
Europe, fundamentally reshaping our understanding of past settlement patterns and landscape use.

However, the interpretation of LiDAR-derived imagery for archaeological purposes remains
overwhelmingly a manual, expert-driven process. Trained archaeologists visually inspect
visualisations tile by tile, identifying and digitising features of interest. At national scales,
this approach is unsustainable: Scotland alone contains approximately 80,000~km$^{2}$ of
archaeologically sensitive terrain, the majority of which has never been systematically
surveyed \citep{banaszek2018only}. The bottleneck is not data acquisition---LiDAR coverage is
expanding rapidly through national programmes---but interpretation capacity. Manual
interpretation introduces further limitations: inter-observer variability, fatigue-related
omissions, and inconsistency in feature classification across different surveyors and projects.

Deep learning, and specifically convolutional neural network (CNN)-based object detection,
offers a promising solution to this interpretation bottleneck. Modern object detection
architectures can learn to recognise complex visual patterns directly from labelled examples,
potentially matching or exceeding human performance on well-defined visual recognition tasks
while operating at computational speed. In archaeology, where target features exhibit
characteristic morphological signatures in LiDAR visualisations, object detection models trained
on annotated examples could automate the first-pass screening of large areas, allowing human
experts to focus on verification and ambiguous cases---a paradigm that has proven highly
effective in medical imaging, remote sensing, and industrial inspection \citep{kadhim2023critical}.

Among the available object detection architectures, the You Only Look Once (YOLO) family has
gained particular prominence due to its combination of speed, accuracy, and accessibility.
First introduced by \citet{redmon2016you}, YOLO formulates object detection as a single
regression problem, predicting bounding boxes and class probabilities directly from image
pixels in one forward pass. Successive iterations---from YOLOv5 through YOLOv8, YOLOv11, and
the most recent YOLOv12---have incorporated architectural advances including anchor-free
detection heads, advanced feature pyramid networks, and attention mechanisms that improve
performance on small and densely clustered objects \citep{jocher2023ultralytics}. These
characteristics make YOLO particularly well-suited to archaeological detection tasks, where
targets are often small relative to image dimensions and occur in clusters of varying density.

This paper presents a systematic evaluation of multi-version YOLO fine-tuning for
archaeological site detection in LiDAR-derived hillshade imagery, using the Isle of Arran
(Scotland) as a test case. We introduce a consensus-based inference approach in which
multiple random seeds of the same model architecture vote on the presence of each detected
feature, requiring agreement from at least two of three seeds before a detection is
retained. The study addresses three core questions: (1)~How effectively do modern YOLO
variants (v8, v11, v12) detect archaeological features in hillshade imagery when fine-tuned
with limited labelled data? (2)~Does multi-seed consensus voting improve the reliability of
detections compared to single-model predictions? (3)~What are the practical implications of
model version choice and consensus threshold for archaeological prospection workflows?

% ======================================================================
\section{Related Work}
% ======================================================================

\subsection{LiDAR Visualisation for Archaeological Prospection}

The application of LiDAR to archaeology began in earnest in the early 2000s, with pioneering
studies demonstrating its ability to penetrate forest canopy and reveal ground-surface
topography at sub-metre resolution \citep{crutchley2010light}. The key innovation that
unlocked the archaeological potential of LiDAR was the development of specialised
visualisation techniques that emphasise subtle topographic features while suppressing
large-scale terrain variation.

\citet{hesse2010lidar} introduced the Local Relief Model (LRM), which decomposes a DTM into
local (small-scale) and regional (large-scale) components through low-pass filtering,
retaining only the local component to highlight archaeological features. This approach
proved particularly effective for detecting low-relief earthworks such as field boundaries,
trackways, and building platforms that are otherwise invisible in standard hillshade
renderings. \citet{kokalj2011application} proposed the Sky-View Factor (SVF), a
visualisation method borrowed from atmospheric sciences that computes the proportion of
visible sky from each point on a surface. SVF exaggerates convex features (banks, cairns,
walls) while suppressing concave features (pits, ditches), producing visually intuitive
renderings that non-specialists can interpret. Subsequent research has extended this
toolkit to include multiscale hillshading, openness, positive and negative openness, and
principal component analysis of multiple visualisation layers \citep{doneus2008archaeological,
devereux2008visualisation, mayoral2020automated}.

These visualisation techniques have been instrumental in the discovery of archaeological
sites at landscape scale. In Scotland, the National Lidar Programme has acquired data
covering substantial portions of the country, and projects such as the Rapid Archaeological
Mapping Programme (RAMP) have demonstrated the potential of LiDAR-based survey for
identifying previously unrecorded sites in upland environments \citep{banaszek2018only}.
Nevertheless, the interpretation process remains overwhelmingly manual, creating an urgent
need for automated detection methods that can keep pace with the growing volume of
available LiDAR data.

\subsection{Deep Learning for Archaeological Detection}

The application of deep learning to archaeological remote sensing has expanded rapidly since
approximately 2018, driven by advances in computer vision, the availability of pre-trained
models, and the growing volume of labelled archaeological data. Two comprehensive reviews
provide a systematic overview of this emerging field.

\citet{argyrou2022review} reviewed 135 studies applying artificial intelligence and remote
sensing to archaeology, documenting a clear shift from traditional machine learning (random
forests, support vector machines) towards deep learning architectures after 2018. They
identified object detection as one of the most promising application areas, with
convolutional neural networks achieving superior performance on tasks ranging from burial
mound detection in satellite imagery to structure identification in drone-based
orthophotography. However, they noted that the majority of studies remained at the proof-of-concept
stage, applied to small geographic areas with limited ground-truth data, and that
reproducibility and generalisability across different landscape contexts remained
significant challenges.

\citet{kadhim2023critical} provided a more technically focused review of deep learning
approaches in archaeology, cataloguing the architectures, data sources, and evaluation
metrics employed across 85 studies. They found that YOLO variants were among the most
frequently used architectures, alongside Faster R-CNN and U-Net, with applications
spanning site detection, artefact classification, and inscription recognition. Key
limitations identified included the scarcity of large, publicly available annotated
datasets; the computational cost of training; and the difficulty of transferring models
between regions with different archaeological signatures and landscape characteristics.

\subsection{Object Detection in LiDAR-Derived Imagery}

Several recent studies have specifically addressed the challenge of detecting archaeological
features in LiDAR-derived visualisations using deep learning, and these are most directly
relevant to the present work.

\citet{canedo2023uncovering} applied a data-centric AI approach to the detection of
archaeological sites in airborne LiDAR data from the Iberian Peninsula. Rather than
focusing on architectural innovation, they systematically investigated the impact of data
quality, annotation strategy, and training set composition on detection performance. Their
results demonstrated that careful data curation---including the selection of appropriate
visualisation types, consistent annotation protocols, and balanced class representation---
could yield performance improvements comparable to or exceeding those from architectural
changes. This finding underscored a critical insight for archaeological applications: in a
domain where labelled data is scarce and expensive to produce, data-centric strategies may
offer higher return on investment than model-centric ones.

Building on this work, \citet{canedo2025automated} developed a deep multimodal segmentation
approach for the automated detection of hillforts in remote sensing imagery. Their method
integrated multiple visualisation modalities---hillshade, slope, and local relief---as
input channels to a segmentation network, enabling the model to learn complementary
features from different representations of the same terrain. The multimodal approach
outperformed single-modality baselines, particularly for features with ambiguous
morphological signatures, demonstrating the value of sensor and visualisation fusion for
archaeological detection.

\citet{marccal2024evaluating} provided the most direct precedent for the present study by
comparing Faster R-CNN and YOLOv8 for the detection of megalithic monuments in satellite
imagery. Their results showed that YOLOv8 achieved competitive accuracy with significantly
faster inference speed compared to Faster R-CNN, reinforcing YOLO's suitability for
large-scale archaeological survey applications where processing speed is a practical
constraint. However, their study was limited to a single YOLO version and did not explore
the impact of random seed variation or consensus-based approaches.

\subsection{The Multi-Seed Consensus Gap}

A notable gap in the existing literature is the absence of systematic investigation into
model reliability and prediction consistency. Deep learning models are inherently stochastic:
different random initialisations (seeds) can lead to different learned weights, and
consequently to different predictions on the same input. In high-stakes applications such
as medical diagnosis and autonomous driving, ensemble and consensus techniques are
routinely employed to mitigate this variability. In archaeology, however, the standard
practice remains to train a single model and report its performance, without accounting for
the variability introduced by the training process.

This gap has practical consequences. An archaeological survey that relies on a single
model's predictions risks both false positives (features incorrectly identified that
require unnecessary ground verification) and false negatives (genuine features missed
entirely). A consensus approach, requiring agreement across independently trained models
before a detection is reported, offers a principled mechanism for controlling this
trade-off. To our knowledge, no prior study has systematically evaluated multi-seed
consensus voting for archaeological object detection, nor compared the behaviour of
different YOLO generations under this paradigm.

% ======================================================================
\section{Research Objectives}
% ======================================================================

The overarching aim of this study is to develop and evaluate a reproducible, multi-version
deep learning pipeline for archaeological site detection in LiDAR-derived hillshade
imagery, with a specific focus on model reliability through consensus-based inference.
The study is conducted on the Isle of Arran, Scotland---a region with a rich archaeological
record spanning the Neolithic to the post-medieval period, and for which high-resolution
LiDAR-derived visualisation data are available.

The specific objectives are:

\begin{enumerate}[leftmargin=*]

\item \textbf{To fine-tune and compare three generations of YOLO nano-scale architectures}
(YOLOv8n, YOLO11n, and YOLO12n) for the detection of three archaeological feature classes---
roundhouses, shielings, and small cairns---in LiDAR-derived hillshade imagery. Each model
version is trained under three different random seeds (0, 1, 2) to characterise the
variability introduced by stochastic initialisation, producing nine independently trained
models in total.

\item \textbf{To implement and evaluate a multi-seed consensus voting mechanism} in which
detections from the three seeds of each model version are compared, and only those
predictions for which at least two of three seeds agree (Intersection-over-Union $\geq$ 0.5,
same predicted class) are retained. This approach provides a quantitative estimate of
prediction reliability that goes beyond single-model confidence scores.

\item \textbf{To compare single-model, per-version consensus, and cross-version ensemble
performance} using standard COCO-style mean Average Precision at 0.5 IoU threshold
(mAP@0.5), supplemented by per-class precision and recall metrics, in order to characterise
the trade-offs between detection sensitivity and false positive rate under different
inference strategies.

\item \textbf{To establish a reproducible, open-source pipeline}---encompassing data
preparation, multi-seed training, consensus inference, and quantitative evaluation---that
can be applied to other archaeological contexts and LiDAR datasets with minimal adaptation
\citep{yingqingyang2026arran}.

\end{enumerate}

By pursuing these objectives, this study contributes to the growing body of evidence on
deep learning applications in archaeology while addressing the critical, and currently
underexplored, issue of prediction reliability through consensus mechanisms. The results
provide practical guidance for archaeological practitioners considering the adoption of
automated detection methods, and for researchers developing the next generation of
archaeological prospection tools.

% ======================================================================
\section{Study Area and Data}
% ======================================================================

The Isle of Arran (55.58$^{\circ}$N, 5.21$^{\circ}$W) is located in the Firth of Clyde,
off the west coast of Scotland. Often described as ``Scotland in miniature'' due to its
diverse geology and landforms, Arran contains a concentration of archaeological sites
spanning from the Neolithic to the post-medieval period, including chambered cairns,
standing stones, roundhouses, shielings, and extensive field systems \citep{barber1997arran}.
The island has been the subject of intensive archaeological survey since the 1970s, and its
archaeological record is relatively well-documented, making it an ideal test case for
evaluating automated detection methods against reliable ground truth.

For this study, we use LiDAR-derived hillshade imagery of Arran processed with
multidirectional illumination. The dataset comprises 424 image tiles of 500 $\times$
500~pixels, manually annotated by archaeological specialists with bounding boxes for three
feature classes: roundhouses (prehistoric circular stone-built dwellings, typically 8--15~m
in diameter), shielings (seasonal upland pastoral structures, typically sub-rectangular,
5--10~m across), and small cairns (clearance or burial cairns, 2--6~m in diameter, often
appearing as subtle mounds in the hillshade). The dataset was split into 364 training
images (85.8\%) and 60 validation images (14.2\%), with annotations in YOLO normalised
format. Instance counts per class in the training set were 260 (roundhouse), 419 (shieling),
and 489 (small cairn), reflecting the relative abundance of these feature types in the
Arran landscape.

% ======================================================================
\bibliographystyle{unsrtnat}
\bibliography{references}
\end{document}
